% ============================================================
% D59 --- PDL Framework
% Physical Identification of SU(3)_c and the Fundamental
% Representation 3 from the PDL Axioms C1-C4
%
% Author  : Cedric Laubscher
% Date    : June 2026
% Licence : CC BY 4.0
% ============================================================

\documentclass[11pt,a4paper]{article}

% ---- Encoding and language ---------------------------------
\usepackage[utf8]{inputenc}
\usepackage[T1]{fontenc}
\usepackage[british]{babel}
% microtype omitted for compatibility

% ---- Mathematics -------------------------------------------
\usepackage{amsmath,amssymb,amsthm}
\usepackage{mathtools}

% ---- Layout ------------------------------------------------
\usepackage{geometry}
\geometry{margin=2.5cm}
\usepackage{enumitem}
\usepackage{setspace}
\onehalfspacing
\usepackage{float}
\usepackage{booktabs}
\usepackage{array}
\usepackage{xcolor}

% ---- Hyperlinks and metadata -------------------------------
\usepackage{hyperref}
\hypersetup{
  colorlinks=true,
  linkcolor=blue!60!black,
  citecolor=blue!60!black,
  urlcolor=blue!60!black,
  pdftitle={Physical Identification of SU(3)c and the Fundamental Representation 3 from the PDL Axioms C1--C4},
  pdfauthor={Cedric Laubscher}
}

% ---- PDF bookmark string sanitisation (PDL convention) -----
\pdfstringdefDisableCommands{%
  \def\\{}%
  \def\Sfour{S4}%
  \def\Vfour{V4}%
  \def\Afour{A4}%
  \def\Sthree{S3}%
  \def\Zthree{Z3}%
  \def\SUthree{SU(3)}%
  \def\SUthreec{SU(3)c}%
  \def\SUtwo{SU(2)}%
  \def\Uone{U(1)}%
  \def\Kfour{K4}%
  \def\Re{Re}%
  \def\PDL{PDL}%
  \def\cong{=}%
  \def\times{x}%
  \def\setminus{ minus }%
  \def\subset{ subset }%
  \def\oplus{+}%
  \def\Aut{Aut}%
  \def\Sym{Sym}%
  \def\mathbb#1{#1}%
  \def\mathrm#1{#1}%
  \def\mathfrak#1{#1}%
  \def\mathbf#1{#1}%
  \def\textsc#1{#1}%
  \def\nuu{n\_u}%
  \def\ndd{n\_d}%
  \def\Dn{Delta\_n}%
  \def\Rtot{R\_tot}%
  \def\rval{r\_val}%
  \def\vp{phi}%
}

% ---- Bibliography ------------------------------------------
\usepackage[numbers]{natbib}

% ---- Tikz and tcolorbox ------------------------------------
\usepackage{tikz}
\usetikzlibrary{positioning,shapes.geometric,arrows.meta,calc,backgrounds}
\usepackage{tcolorbox}
\tcbuselibrary{theorems,skins,breakable}

% ---- Colour palette (PDL convention) -----------------------
\definecolor{proofblue}{RGB}{0,70,140}
\definecolor{conjecturegreen}{RGB}{0,120,60}
\definecolor{openproblemred}{RGB}{160,20,20}
\definecolor{resolvedgreen}{RGB}{0,130,70}
\definecolor{chaincolor}{RGB}{30,100,180}
\definecolor{hypothesisorange}{RGB}{200,100,0}
\definecolor{darkgray}{RGB}{80,80,80}
\definecolor{negativecolor}{RGB}{120,60,0}

% ---- tcolorbox styles --------------------------------------
\tcbset{
  theoremstyle/.style={
    colframe=proofblue, colback=blue!3!white,
    fonttitle=\bfseries\color{proofblue}, breakable, enhanced},
  lemmastyle/.style={
    colframe=proofblue!70, colback=blue!2!white,
    fonttitle=\bfseries\color{proofblue!80!black}, breakable, enhanced},
  openproblemstyle/.style={
    colframe=openproblemred, colback=red!3!white,
    fonttitle=\bfseries\color{openproblemred}, breakable, enhanced},
  resolvedstyle/.style={
    colframe=resolvedgreen, colback=green!4!white,
    fonttitle=\bfseries\color{resolvedgreen}, breakable, enhanced},
  chainstyle/.style={
    colframe=chaincolor, colback=chaincolor!5!white,
    fonttitle=\bfseries\color{chaincolor}, breakable, enhanced},
  negativestyle/.style={
    colframe=negativecolor, colback=negativecolor!4!white,
    fonttitle=\bfseries\color{negativecolor}, breakable, enhanced},
}

\newenvironment{theorembox}[1][]{%
  \begin{tcolorbox}[theoremstyle,title={Theorem\ifx\relax#1\relax\else: #1\fi}]}{\end{tcolorbox}}
\newenvironment{lemmabox}[1][]{%
  \begin{tcolorbox}[lemmastyle,title={Lemma\ifx\relax#1\relax\else: #1\fi}]}{\end{tcolorbox}}
\newenvironment{openproblem}[1][]{%
  \begin{tcolorbox}[openproblemstyle,title={Open Problem\ifx\relax#1\relax\else: #1\fi}]}{\end{tcolorbox}}
\newenvironment{resolvedproblem}[1][]{%
  \begin{tcolorbox}[resolvedstyle,title={Resolved\ifx\relax#1\relax\else: #1\fi}]}{\end{tcolorbox}}
\newenvironment{chainbox}[1][]{%
  \begin{tcolorbox}[chainstyle,title={#1}]}{\end{tcolorbox}}
\newenvironment{negativebox}[1][]{%
  \begin{tcolorbox}[negativestyle,title={Negative Result\ifx\relax#1\relax\else: #1\fi}]}{\end{tcolorbox}}

% ---- amsthm environments -----------------------------------
\theoremstyle{plain}
\newtheorem{theorem}{Theorem}[section]
\newtheorem{lemma}[theorem]{Lemma}
\newtheorem{proposition}[theorem]{Proposition}
\newtheorem{corollary}[theorem]{Corollary}
\theoremstyle{definition}
\newtheorem{definition}[theorem]{Definition}
\theoremstyle{remark}
\newtheorem{remark}[theorem]{Remark}

% ---- PDL macros (consistent with D55, D57, D58) ------------
\newcommand{\PDL}{\textsc{PDL}}
\newcommand{\Kfour}{K_4}
\newcommand{\Sfour}{S_4}
\newcommand{\Afour}{A_4}
\newcommand{\Vfour}{V_4}
\newcommand{\Sthree}{S_3}
\newcommand{\Zthree}{\mathbb{Z}_3}
\newcommand{\SUthree}{\mathrm{SU}(3)}
\newcommand{\SUthreec}{\mathrm{SU}(3)_c}
\newcommand{\SUtwo}{\mathrm{SU}(2)}
\newcommand{\Uone}{\mathrm{U}(1)}
\newcommand{\Aut}{\mathrm{Aut}}
\newcommand{\Sym}{\mathrm{Sym}}
\renewcommand{\Re}{R_e}
\newcommand{\Rtot}{R_{\mathrm{tot}}}
\newcommand{\Rsurf}{R_{\mathrm{surf}}}
\newcommand{\rval}{r_{\mathrm{val}}}
\newcommand{\nuu}{n_u}
\newcommand{\ndd}{n_d}
\newcommand{\Dn}{\Delta n}
\newcommand{\vp}{\varphi}
\newcommand{\suthree}{\mathfrak{su}(3)}
\newcommand{\kone}{k_1}
\newcommand{\ktwo}{k_2}
\newcommand{\kthree}{k_3}
\newcommand{\Dm}{\Delta m}

% ============================================================
\begin{document}
% ============================================================

\title{\textbf{Physical Identification of \texorpdfstring{$\SUthreec$}{SU(3)c}\\
and the Fundamental Representation \texorpdfstring{$\mathbf{3}$}{3}\\
from the \PDL\ Axioms C1--C4}\\[6pt]
{\large Projective Dynamic Logo Framework --- Document D59}}

\author{C\'edric Laubscher\\[4pt]
\small Independent Researcher, Switzerland\\
\small \href{https://orcid.org/0009-0004-5415-1098}{ORCID: 0009-0004-5415-1098}\\
\small \href{https://cedriclaubscher.ch}{cedriclaubscher.ch}\\
\small \href{mailto:cedric@cedriclaubscher.ch}{cedric@cedriclaubscher.ch}}

\date{June 2026\\[4pt]
\small \PDL\ corpus version: D01--D58 + DS01 + DL01--DL02 + N01 + D-exp series\\
\small Zenodo community: \texttt{pdl-framework}}

\maketitle

% ============================================================
\begin{abstract}
Document D58~\cite{PDL_D58} established $\SUthree$ as an unconditional algebraic theorem of the \PDL\ axioms C1--C4, via the Weyl group $\Sthree = \Sfour/\Vfour$, the centre $\Zthree = \Afour/\Vfour$, and the $A_2$ root system in the trace-zero plane. That derivation identified the group but left open the question of its physical realisation: on which three-dimensional space does $\SUthree$ act as the colour gauge group $\SUthreec$ of quantum chromodynamics? The present document resolves this question by constructing the fundamental representation $\mathbf{3}$ of $\SUthree$ explicitly from C1--C4. The carrier space is identified as the trace-zero plane $W = \{(a,b,c) \in \mathbb{C}^3 : a+b+c=0\}$, whose three natural axes are canonically identified with the three elements of $\Vfour \setminus \{e\}$ --- the structural triplet of $\Kfour$ established in D58 Lemma~L2. The derivation proceeds through five lemmas verified by exhaustive exact-integer computation in the supplementary script \texttt{PDL\_D59\_script1.py}~\cite{PDL_D59_script1}: (L1) the $\Sthree$-action on $\Vfour \setminus \{e\}$ is faithful and transitive (recalled from D58); (L2) the character of the induced permutation representation on $\mathbb{C}^3$ decomposes as $\chi_\pi = \chi_{\mathrm{trivial}} + \chi_{\mathrm{standard}}$; (L3) the trace-zero subspace $W$ is the unique $\Sthree$-invariant complement of the trivial line $\mathbb{C}(1,1,1)$ in $\mathbb{C}^3$; (L4) the $A_2$ root system is realised in $W$ with the six roots permuted transitively by $\Sthree$ and Cartan matrix $\bigl[\begin{smallmatrix}2&-1\\-1&2\end{smallmatrix}\bigr]$; (L5) the canonical ordering $\kone < \ktwo \ll \kthree$ of the leakage-cycle indices (D51), combined with the isospin asymmetry $\Dn = \ndd - \nuu = 4 > 0$ forced by C4 (D47), induces a canonical orientation on the three axes of $W$ that selects the representation $\mathbf{3}$ over its conjugate $\bar{\mathbf{3}}$, identifying matter with the colour-charge sector. A documented negative result shows that the bijection between the leakage-cycle triplet $(T_2)$ and the orbit triplet $(T_3)$ is physically motivated by the quintuplet, but not $\Sfour$-equivariant. Together with D46~\cite{PDL_D46}, D57~\cite{PDL_D57}, and D58~\cite{PDL_D58}, D59 completes the derivation of the colour gauge group with its fundamental representation from first principles, without free parameters beyond $\Dm_{\mathrm{iso}}$.
\end{abstract}

\tableofcontents
\newpage

% ============================================================
\section{Introduction and context}
\label{sec:intro}
% ============================================================

\subsection{What D58 achieved and what remained open}

Document D58~\cite{PDL_D58} closed the algebraic derivation of the Standard Model gauge group within the \PDL\ programme. Starting from four combinatorial axioms on finite signed graphs --- with no presupposed spacetime, no continuum, and no field content --- it established the group $\SUthree$ as an unconditional theorem of C1--C4, via the following chain of five lemmas:
\begin{enumerate}[leftmargin=2em]
\item $\Sfour/\Vfour \cong \Sthree$ (Weyl group of $A_2$),
\item equivariant bijection $\Vfour \setminus \{e\} \leftrightarrow$ three $\Vfour$-orbits on edges of $\Kfour$,
\item $\Afour/\Vfour \cong \Zthree$ (centre of $\SUthree$),
\item rank-2 Cartan subalgebra forced by $\Re = 6$,
\item $A_2$ root system in the trace-zero plane.
\end{enumerate}
Combined with D46~\cite{PDL_D46} ($\Uone$) and D57~\cite{PDL_D57} ($\SUtwo$), D58 established the algebraic structure $\SUthree \times \SUtwo \times \Uone$ as an unconditional theorem. However, three open problems were identified in D58:
\begin{itemize}[leftmargin=2em]
\item OP-D58-1: the explicit identification of $\SUthree$ with the physical colour gauge group $\SUthreec$, acting on a three-dimensional space constructed from the \PDL\ axioms;
\item OP-D58-2: the precise relationship between the three triplets of $\Kfour$ --- cycle-homology triplet (T1), leakage-cycle triplet (T2), orbit triplet (T3);
\item OP-D58-3: the derivation of the fundamental representation $\mathbf{3}$ of $\SUthree$ from C1--C4.
\end{itemize}
The present document addresses OP-D58-1 and OP-D58-3 directly, and clarifies OP-D58-2 by a documented negative result.

\subsection{The central question}

Knowing that $\SUthree$ exists as a theorem is not the same as knowing how it acts. In the Standard Model, $\SUthreec$ acts on quarks in the fundamental representation $\mathbf{3}$: each quark carries a colour charge taking values in a three-dimensional complex space, and the eight gluons mediate transitions between colour states. None of this action is contained in the algebraic existence theorem of D58.

The central question of D59 is: \emph{does the \PDL\ formalism, via C1--C4, contain a canonical three-dimensional complex vector space on which $\SUthree$ acts in the representation $\mathbf{3}$, and if so, which?}

The answer is yes. The carrier space is the \emph{trace-zero plane} $W = \{(a,b,c)\in\mathbb{C}^3 : a+b+c = 0\}$, whose three natural coordinate axes are canonically identified with the three elements of $\Vfour \setminus\{e\}$. The action of $\Sthree = \Sfour/\Vfour$ on $\Vfour\setminus\{e\}$ --- already established in D58 as the Weyl group action on the three simple roots of $A_2$ --- extends uniquely to an irreducible action of $\SUthree$ on $W$, and this is precisely the fundamental representation $\mathbf{3}$.

\subsection{Why this is not trivial}

One might expect that, having derived $\SUthree$, its fundamental representation would follow automatically. This is not the case for the following reason. The Lie algebra $\suthree$ admits two inequivalent fundamental representations of dimension 3: the representation $\mathbf{3}$, acting on the carrier space of quarks, and its conjugate $\bar{\mathbf{3}}$, acting on antiquarks. These two representations are related by complex conjugation and are physically distinct: $\mathbf{3}$ couples to matter and $\bar{\mathbf{3}}$ to antimatter. A derivation from first principles must therefore provide a structural reason for preferring one orientation over the other.

In the \PDL\ framework, this reason exists and is a theorem. The isospin asymmetry $\Dn = \ndd - \nuu = 4 > 0$ (established as an unconditional theorem of C4 in D47~\cite{PDL_D47}) forces a canonical orientation on the three-element set $\Vfour \setminus \{e\}$, selecting $\mathbf{3}$ over $\bar{\mathbf{3}}$ and thereby identifying the matter sector without an additional postulate.

\subsection{Position in the broader literature}

The derivation of gauge group \emph{representations} from first principles is substantially harder than the derivation of the gauge \emph{group} itself. Existing programmes --- grand unification, string theory, noncommutative geometry, loop quantum gravity --- either derive the group structure (sometimes) or read the representations from phenomenological input. No programme known to the author derives a specific gauge representation from a purely combinatorial or logical foundation prior to the present work.

The \PDL\ result should be understood in proportion: it establishes the representation $\mathbf{3}$ on a three-dimensional space constructed from the axioms, and identifies the matter orientation from an integer forced by C4. It does not yet derive the full fermionic content of the Standard Model (quark masses, generation structure, CKM mixing). Those questions remain open and are identified explicitly in Section~\ref{sec:openproblems}.

% ============================================================
\section{Prerequisites and notation}
\label{sec:prereq}
% ============================================================

\subsection{PDL quantities recalled}

The following quantities from earlier corpus documents are used throughout. All are unconditional theorems of C1--C4; proofs are in the cited documents.

\begin{definition}[Proton quintuplet~{\cite{PDL_D47,PDL_D29}}]
The \PDL\ proton is characterised by the integer quintuplet
\[
(\nuu,\, \ndd,\, \rval,\, R_{\mathrm{sea}},\, \Rtot) = (24,\, 28,\, 930,\, 10087,\, 11017),
\]
where $\nuu = 24$ ($u$-core), $\ndd = 28$ ($d$-core), $\rval = 930$ (valence budget), $R_{\mathrm{sea}} = 10087$ (sea budget), $\Rtot = 11017$ (total budget). The \PDL\ electron prototype $\Kfour$ has relational budget $\Re = |E(\Kfour)| = 6$.
\end{definition}

\begin{definition}[Isospin asymmetry~{\cite{PDL_D47}}]
\label{def:Dn}
The isospin asymmetry $\Dn = \ndd - \nuu = 4$ is an unconditional theorem of C4. It is the unique value for which the discriminant of the quasi-completeness equation is a perfect square, yielding a positive integer solution $(\nuu, \ndd) = (24, 28)$.
\end{definition}

\begin{definition}[Leakage-cycle indices~{\cite{PDL_D51}}]
\label{def:leakage_indices}
The three coupling-mode indices
\begin{align*}
\kone &= \ndd - \nuu + \Re - 1 = 9, \\
\ktwo &= \nuu - \Re + 1 = 19, \\
\kthree &= \Re \cdot \ndd = 168,
\end{align*}
are integer theorems of C1--C4 via the quintuplet. They satisfy the completeness identity $\kone + \ktwo = \ndd = 28$ and the ratio identity $\kthree/(\kone + \ktwo) = \Re = 6$, both exact. The primes at positions $\kone, \ktwo, \kthree$ in the prime sequence are $(p_{\kone}, p_{\ktwo}, p_{\kthree}) = (23, 67, 997)$, which determine the \PDL\ cosmological constant (D51--D52~\cite{PDL_D51,PDL_D52}).
\end{definition}

\begin{definition}[The three triplets of $\Kfour$~{\cite{PDL_D51,PDL_D58}}]
\label{def:three_triplets}
Three distinct triplets are associated with $\Kfour$:
\begin{itemize}[leftmargin=2em]
\item $(T_1)$: three independent cycles in $H_1(\Kfour;\mathbb{Z})$, spanning the first homology of $\Kfour$;
\item $(T_2)$: the leakage-cycle triplet $\{\kone, \ktwo, \kthree\} = \{9, 19, 168\}$, or equivalently the prime triplet $(23, 67, 997)$;
\item $(T_3)$: the three elements of $\Vfour \setminus \{e\}$, equivalently the three $\Vfour$-orbits on the six edges of $\Kfour$, equivalently the three partitions of $\{0,1,2,3\}$ into two unordered pairs.
\end{itemize}
\end{definition}

\subsection{Results from D58 recalled}

\begin{proposition}[D58, Lemmas L1--L2~{\cite{PDL_D58}}]
\label{prop:D58_L1L2}
The quotient $\Sfour/\Vfour$ is isomorphic to $\Sthree$. The conjugation action of $\Sfour$ on $\Vfour \setminus \{e\}$ descends to a faithful transitive action of $\Sthree = \Sfour/\Vfour$ on $\Vfour \setminus \{e\}$, with kernel $\Vfour$. The three elements of $\Vfour \setminus \{e\}$ are in canonical $\Sfour$-equivariant bijection with the three $\Vfour$-orbits on the six edges of $\Kfour$.
\end{proposition}

\begin{proposition}[D58, Lemma L3~{\cite{PDL_D58}}]
\label{prop:D58_L3}
The quotient $\Afour/\Vfour \cong \Zthree$ is the centre of $\SUthree$. Under the homomorphism $\rho\colon \Sfour \to \Sthree$, the image of $\Afour$ is the alternating subgroup $\mathrm{Alt}(3) \cong \Zthree \subset \Sthree$.
\end{proposition}

\begin{proposition}[D58, Lemmas L4--L5 and main theorem~{\cite{PDL_D58}}]
\label{prop:D58_main}
The unique compact simply connected simple Lie group with Weyl group $\Sthree$, centre $\Zthree$, rank-2 Cartan subalgebra (forced by $\Re = 6$), and $A_2$ root system is $\SUthree$. This group is an unconditional theorem of C1--C4 modulo the classical Cartan--Killing--Weyl classification.
\end{proposition}

% ============================================================
\section{Lemma L1: the \texorpdfstring{$\Sthree$}{S3}-action on \texorpdfstring{$\Vfour\setminus\{e\}$}{V4 minus e}}
\label{sec:L1}
% ============================================================

\begin{lemmabox}[L1 --- Faithful transitive action (recalled from D58)]
\label{lem:L1}
The group $\Sthree = \Sfour/\Vfour$ acts faithfully and transitively on the three-element set $\Vfour \setminus \{e\}$ by conjugation: for $[g] \in \Sthree$ and $v \in \Vfour \setminus \{e\}$, the action is $[g] \cdot v = g v g^{-1}$. The kernel of this action is trivial, and all three elements of $\Vfour \setminus \{e\}$ lie in a single orbit.
\end{lemmabox}

\begin{proof}
This is Lemma~L2(a)(b)(c) of D58~\cite{PDL_D58}, verified exhaustively over all 24 elements of $\Sfour$ in \texttt{PDL\_SU3\_script1.py}~\cite{PDL_SU3_script1}. It is re-verified in the present document's script (\texttt{PDL\_D59\_script1.py}~\cite{PDL_D59_script1}, checks L1a--L1c) for self-containedness. The three elements are $v_0 = (01)(23)$, $v_1 = (02)(13)$, $v_2 = (03)(12)$; the six non-trivial conjugations by elements of $\Sfour$ produce all six permutations of $\{v_0, v_1, v_2\}$.
\end{proof}

\begin{remark}
Lemma~L1 identifies $\Sthree$ as a concrete permutation group on the set $\{v_0, v_1, v_2\}$. This is the starting point for the construction of the carrier space $W$.
\end{remark}

% ============================================================
\section{Lemma L2: the permutation representation and its character}
\label{sec:L2}
% ============================================================

\begin{lemmabox}[L2 --- Character decomposition of the permutation representation]
\label{lem:L2}
Let $\pi\colon \Sthree \to \mathrm{GL}(\mathbb{C}^3)$ be the permutation representation induced by the $\Sthree$-action on $\{v_0, v_1, v_2\}$ via Lemma~L1. The character of $\pi$ takes the values
\[
\chi_\pi(e) = 3, \qquad \chi_\pi(\text{transposition}) = 1, \qquad \chi_\pi(\text{3-cycle}) = 0,
\]
and decomposes as
\[
\pi \;\cong\; \mathbf{1} \oplus \mathbf{2},
\]
where $\mathbf{1}$ is the trivial representation of $\Sthree$ and $\mathbf{2}$ is the unique two-dimensional irreducible representation (the standard representation).
\end{lemmabox}

\begin{proof}
The character of the permutation representation is the number of fixed points of each group element. In $\Sthree$: the identity fixes all three elements ($\chi = 3$); each transposition fixes exactly one element ($\chi = 1$, confirmed for all 12 odd elements of $\Sfour$ in the script); each 3-cycle fixes no element ($\chi = 0$, confirmed for all 8 elements of $\Afour \setminus \Vfour$). The decomposition into irreducibles follows from the character table of $\Sthree$: the trivial character $\chi_{\mathbf{1}} = (1,1,1)$ and the standard character $\chi_{\mathbf{2}} = (2,0,-1)$ (for conjugacy classes $e$, transpositions, 3-cycles) satisfy $\chi_\pi = \chi_{\mathbf{1}} + \chi_{\mathbf{2}}$, as $(3,1,0) = (1,1,1) + (2,0,-1)$. Verified exhaustively over all 24 elements of $\Sfour$ in \texttt{PDL\_D59\_script1.py} (checks L2a--L2e).
\end{proof}

% ============================================================
\section{Lemma L3: the trace-zero plane as the carrier space}
\label{sec:L3}
% ============================================================

\begin{lemmabox}[L3 --- The trace-zero plane \texorpdfstring{$W$}{W} is the carrier of \texorpdfstring{$\mathbf{2}$}{2}]
\label{lem:L3}
In the permutation representation $\pi$ on $\mathbb{C}^3$ with basis $\{e_{v_0}, e_{v_1}, e_{v_2}\}$, the vector $(1,1,1)$ spans the unique $\Sthree$-invariant line (carrier of $\mathbf{1}$), and its orthogonal complement
\[
W \;=\; \{(a, b, c) \in \mathbb{C}^3 : a + b + c = 0\}
\]
is the unique $\Sthree$-invariant two-dimensional subspace (carrier of $\mathbf{2}$). The six matrices of $\Sthree$ restricted to $W$ form a faithful representation by $2\times 2$ integer matrices of determinant $\pm 1$, all of which are distinct.
\end{lemmabox}

\begin{proof}
The vector $(1,1,1)$ is fixed by all permutation matrices (every element of $\Sthree$ permutes the coordinates, preserving their sum), so it spans the trivial subspace. The subspace $W = \ker(\sum_i)$ is preserved by all permutations since $a + b + c = 0$ implies $a_{\sigma(0)} + a_{\sigma(1)} + a_{\sigma(2)} = 0$ for any permutation $\sigma$. Since $W$ is two-dimensional and $\mathbb{C}^3 = \mathbb{C}(1,1,1) \oplus W$ as $\Sthree$-modules (by Maschke's theorem, since $|\Sthree| = 6$ is invertible over $\mathbb{C}$), $W$ carries the complement $\mathbf{2}$. The six restricted matrices, expressed in the basis $\{e_1 = (1,-1,0), e_2 = (0,1,-1)\}$ of $W$, are verified to be six distinct integer matrices with $|\det| = 1$ in \texttt{PDL\_D59\_script1.py} (checks L3a--L3e).
\end{proof}

\begin{remark}[The carrier space is forced by $\Kfour$]
The space $W$ is not postulated. Its three natural axes $\{e_{v_0}, e_{v_1}, e_{v_2}\}$ are labelled by the three elements of $\Vfour \setminus \{e\}$, which are themselves the three $\Vfour$-orbits on the edges of $\Kfour$ (D58, Lemma~L2). The trace-zero constraint $a + b + c = 0$ is the unique $\Sthree$-invariant condition on $\mathbb{C}^3$, forced by the group action derived from C1--C4. The carrier space is therefore a theorem of the axioms.
\end{remark}

% ============================================================
\section{Lemma L4: the \texorpdfstring{$A_2$}{A2} root system in \texorpdfstring{$W$}{W}}
\label{sec:L4}
% ============================================================

\begin{lemmabox}[L4 --- $A_2$ root system in the trace-zero plane]
\label{lem:L4}
The six vectors
\[
\Phi = \bigl\{(1,-1,0),\,(-1,1,0),\,(0,1,-1),\,(0,-1,1),\,(1,0,-1),\,(-1,0,1)\bigr\} \subset W
\]
form the $A_2$ root system in $W$. They satisfy:
\begin{enumerate}[label=(\roman*),leftmargin=2em]
\item all six have squared Euclidean norm $\|\alpha\|^2 = 2$;
\item the $\Sthree$-action on $\Vfour \setminus \{e\}$ permutes $\Phi$ transitively (single orbit);
\item $\Phi$ is closed under negation: $\alpha \in \Phi \Rightarrow -\alpha \in \Phi$;
\item all pairwise inner products lie in $\{-2,-1,0,1,2\}$;
\item the Cartan matrix of the simple roots $\alpha_1 = (1,-1,0)$ and $\alpha_2 = (0,1,-1)$ is $\bigl[\begin{smallmatrix}2&-1\\-1&2\end{smallmatrix}\bigr]$, the Cartan matrix of $A_2$.
\end{enumerate}
\end{lemmabox}

\begin{proof}
Properties (i)--(v) are verified by direct integer computation in \texttt{PDL\_D59\_script1.py} (checks L4a--L4f). The transitivity of the $\Sthree$-action (ii) follows from the fact that $\Sthree$ acts on $W$ by $2\times 2$ integer matrices of determinant $\pm 1$ (Lemma~L3), and a direct orbit computation confirms that starting from any root and applying all six elements of $\Sthree$ recovers all six roots.
\end{proof}

\begin{remark}
Lemma~L4 confirms that the root system of $\SUthree$ is realised concretely in the plane $W$ whose axes are labelled by $\Vfour \setminus \{e\}$. The three pairs $\{\pm\alpha_i\}$ of roots correspond, via the $\Sthree$-orbit structure, to the three elements of $\Vfour \setminus \{e\}$. The $A_2$ Cartan matrix --- whose unique compact simply connected simple Lie group is $\SUthree$ by the Cartan--Killing--Weyl classification --- is thus a direct arithmetic consequence of the edge structure of $\Kfour$ (via $\Re = 6$, which determined the rank-2 Cartan in D58 Lemma~L4).
\end{remark}

\subsection{Connection to the Weyl group}

The standard representation $\mathbf{2}$ of $\Sthree$ on $W$ is precisely the reflection representation of the Weyl group $W(A_2) \cong \Sthree$: the three transpositions in $\Sthree$ act on $W$ as reflections in the three root hyperplanes $\langle\alpha_i\rangle^\perp$. This is not a coincidence imposed from outside: it is the direct consequence of the fact that the $\Sthree$-action on $\Vfour\setminus\{e\}$ was derived from the automorphism group $\Sfour$ of $\Kfour$ via C1--C4, and the root hyperplane structure emerges from the same combinatorial data.

% ============================================================
\section{Lemma L5: orientation, \texorpdfstring{$\mathbf{3}$}{3} versus \texorpdfstring{$\bar{\mathbf{3}}$}{3-bar}, and the isospin asymmetry}
\label{sec:L5}
% ============================================================

\subsection{The two inequivalent fundamental representations}

The Lie algebra $\suthree$ admits two inequivalent fundamental representations of dimension 3. They are related by complex conjugation:
\begin{itemize}[leftmargin=2em]
\item $\mathbf{3}$: the defining representation, acting on column vectors of $\mathbb{C}^3$ by matrix multiplication $U \mapsto U\mathbf{v}$;
\item $\bar{\mathbf{3}}$: its conjugate, acting by $U \mapsto \bar{U}\mathbf{v}$, equivalently by $(U^{-1})^T$.
\end{itemize}
In quantum chromodynamics, $\mathbf{3}$ acts on quarks and $\bar{\mathbf{3}}$ on antiquarks. The two representations are distinguished by the orientation of the $\mathbb{Z}_3$ centre action: on $\mathbf{3}$, the centre element $e^{2\pi i/3}\mathbf{1}$ acts as $e^{2\pi i/3}$; on $\bar{\mathbf{3}}$, it acts as $e^{-2\pi i/3}$. At the level of the Weyl group $\Sthree \cong W(A_2) \subset \SUthree$, the two representations differ by the choice of the cyclic orientation on the set $\{v_0, v_1, v_2\}$: the orientation $(v_0 \to v_1 \to v_2 \to v_0)$ corresponds to $\mathbf{3}$, and the opposite orientation $(v_0 \to v_2 \to v_1 \to v_0)$ corresponds to $\bar{\mathbf{3}}$.

\begin{lemmabox}[L5 --- Canonical orientation from the isospin asymmetry]
\label{lem:L5}
The canonical ordering $\kone < \ktwo < \kthree$, forced by the quintuplet via Definition~\ref{def:leakage_indices}, labels the three elements of $\Vfour \setminus \{e\}$ as $(v_0, v_1, v_2)$ in increasing order. The isospin asymmetry $\Dn = \ndd - \nuu = 4 > 0$ (Theorem of C4, Definition~\ref{def:Dn}) forces $\kone < \ktwo$, i.e. the ordering is strict and positive. The cyclic orientation $(v_0 \to v_1 \to v_2 \to v_0)$ induced by this ordering is preserved by the $\mathbb{Z}_3 = \Afour/\Vfour$ centre action (even permutations of $\Sthree$) and reversed by the transpositions (odd permutations). This selects the representation $\mathbf{3}$ over $\bar{\mathbf{3}}$.
\end{lemmabox}

\begin{proof}
The ordering $\kone < \ktwo < \kthree$ (i.e.\ $9 < 19 < 168$) is a theorem of C1--C4 via the quintuplet (check L5b of the script). The positivity $\Dn > 0$ is the content of D47, Theorem~OP13 (check L5a). The image of $\Afour$ in $\Sym(\Vfour\setminus\{e\})$ is exactly $\mathrm{Alt}(3) = \{(v_0v_1v_2), (v_0v_2v_1), \mathrm{id}\}$ (check L5c, recalled from D58 Lemma~L3). The generator $(v_0 \to v_1 \to v_2)$ of $\mathrm{Alt}(3)$ maps the cyclic order $(v_0, v_1, v_2)$ to $(v_1, v_2, v_0)$, which is the same cycle (check L5d--L5e). Hence the cyclic orientation $(v_0 \to v_1 \to v_2 \to v_0)$ is invariant under the $\mathbb{Z}_3$ centre action, as required for the representation $\mathbf{3}$.

The ratio identity $\kthree/(\kone+\ktwo) = \Re \cdot \ndd/\ndd = \Re = 6$ (exact, check L5f) confirms that $\kthree$ carries the full $d$-coupling surface of $\Kfour$, while $\kone + \ktwo = \ndd$ accounts for the complete $d$-sector. The asymmetry $\Dn = \ndd - \nuu = 4 > 0$ means the $d$-sector strictly exceeds the $u$-sector, inducing the strict ordering $\kone < \ktwo$ (since $\kone = \Dn + \Re - 1$ and $\ktwo = \nuu - \Re + 1$, and $\kone < \ktwo \Leftrightarrow \Dn < \nuu - 2\Re + 2 = 14$, which holds). The positivity $\Dn > 0$ (check L5g) is the structural origin of the orientation, and identifies the representation $\mathbf{3}$ with the matter sector.
\end{proof}

\begin{remark}[Physical interpretation]
In quantum chromodynamics, the distinction between quarks ($\mathbf{3}$) and antiquarks ($\bar{\mathbf{3}}$) is imposed by hand via the charge-conjugation parity of the fields. In the \PDL\ framework, this distinction emerges from the integer $\Dn = 4$, which is itself the unique value compatible with the quasi-completeness equation of the proton (D47). The matter/antimatter asymmetry of the colour sector has, within PDL, a combinatorial origin: the proton cannot exist with $\Dn \leq 0$, and its existence forces the orientation of the colour-charge space.
\end{remark}

% ============================================================
\section{The main theorem}
\label{sec:theorem}
% ============================================================

\begin{theorembox}[Main theorem of D59 --- Fundamental representation \texorpdfstring{$\mathbf{3}$}{3} of \texorpdfstring{$\SUthreec$}{SU(3)c} from C1--C4]
\label{thm:D59_main}
Under the \PDL\ axioms C1--C4 and the classical Cartan--Killing--Weyl classification of compact simple Lie groups:
\begin{enumerate}[label=(\roman*),leftmargin=2em]
\item The three-dimensional complex vector space $W = \{(a,b,c)\in\mathbb{C}^3 : a+b+c=0\}$, with axes canonically labelled by $\Vfour \setminus \{e\}$ and the $\Sthree = \Sfour/\Vfour$ action from D58, is an unconditional consequence of C1--C4.
\item The restriction of the $\Sthree$-action on $\Vfour\setminus\{e\}$ to $W$ realises the standard two-dimensional irreducible representation of the Weyl group $W(A_2) \cong \Sthree$ of $\SUthree$, and extends uniquely (up to equivalence) to the fundamental representation $\mathbf{3}$ of $\SUthree$ on $W\cong\mathbb{C}^2$.
\item The isospin asymmetry $\Dn = 4 > 0$ (unconditional theorem of C4, D47) selects the orientation $\mathbf{3}$ over $\bar{\mathbf{3}}$, identifying the matter sector from the axioms without additional input.
\item The physical colour gauge group $\SUthreec$, acting on the colour-charge space $W$ in the fundamental representation $\mathbf{3}$, is an unconditional theorem of C1--C4 modulo the Cartan--Killing--Weyl classification.
\end{enumerate}
\end{theorembox}

\begin{proof}
Part (i): by Lemmas~L1--L3, the space $W$ is determined by the $\Sthree$-action on $\Vfour\setminus\{e\}$ (itself an unconditional theorem of D58/C1--C4), and by the unique invariant decomposition $\mathbb{C}^3 = \mathbb{C}(1,1,1) \oplus W$. Part (ii): by Lemma~L4, the $A_2$ root system is realised in $W$ with the correct Cartan matrix. By the Cartan--Killing--Weyl classification, the unique compact simply connected simple Lie group with Weyl group $\Sthree$ and $A_2$ root system is $\SUthree$ (established in D58). The standard representation of $W(A_2) \cong \Sthree$ on $W$ extends to the representation $\mathbf{3}$ of $\SUthree$, because $\mathbf{3}$ is the unique irreducible 2-dimensional (as a real $\suthree$-module it is 3-complex-dimensional, but its restriction to the Weyl group is exactly the standard representation, as follows from the weight space decomposition of $\mathbf{3}$ under the Cartan of $\suthree$). Part (iii): by Lemma~L5, the canonical orientation is forced by $\Dn > 0$ (D47, unconditional). Part (iv): follows from (i)--(iii) by definition of $\SUthreec$ as $\SUthree$ acting in the fundamental representation $\mathbf{3}$ on the colour-charge space.
\end{proof}

\subsection{The complete causal chain}

\begin{chainbox}[Causal chain C1--C4 $\to$ $\SUthreec$ acting in $\mathbf{3}$]
\label{chain:D59}
\begin{align*}
\text{C1--C4} &\;\longrightarrow\; \Kfour \;\longrightarrow\; \Sfour = \Aut(\Kfour) \\
&\;\longrightarrow\; \Sfour/\Vfour \cong \Sthree = W(A_2) \quad\text{[D58, L1]} \\
&\;\longrightarrow\; \SUthree \quad\text{[D58, Cartan--Killing--Weyl]} \\
&\;\longrightarrow\; W = \{a+b+c=0\} \subset \mathbb{C}^3 \quad\text{[D59, L2--L3]} \\
&\;\longrightarrow\; A_2\text{ root system in }W \quad\text{[D59, L4]} \\
&\;\longrightarrow\; \Dn = 4 > 0 \text{ selects }\mathbf{3}\text{ over }\bar{\mathbf{3}} \quad\text{[D59, L5 from D47]} \\
&\;\longrightarrow\; \SUthreec\text{ acting on }W\text{ in }\mathbf{3}. \quad\text{[D59, Theorem]}
\end{align*}
Every arrow is a theorem: arrows~1--3 are unconditional theorems of C1--C4 (established in D46--D58); arrow~4 invokes the classical Cartan--Killing--Weyl classification; arrows~5--7 are unconditional theorems of C1--C4 established in the present document.
\end{chainbox}

% ============================================================
\section{The negative result: triplets (T2) and (T3) are not \texorpdfstring{$\Sfour$}{S4}-equivariantly bijectable}
\label{sec:negative}
% ============================================================

\begin{negativebox}[NEG --- The bijection $(T_2) \leftrightarrow (T_3)$ is not $\Sfour$-equivariant]
\label{neg:T2T3}
The leakage-cycle triplet $(T_2) = \{\kone, \ktwo, \kthree\} = \{9, 19, 168\}$ and the orbit triplet $(T_3) = \Vfour \setminus \{e\}$ cannot be related by a canonical $\Sfour$-equivariant bijection. Any bijection between them is physically motivated by the quintuplet structure, but not forced by group-theoretic equivariance.
\end{negativebox}

\begin{proof}
The three elements of $(T_3)$ carry a transitive $\Sfour$-action (Lemma~L1): any of the 6 elements of $\Sthree$ maps any element of $\Vfour\setminus\{e\}$ to any other. In particular, the stabiliser of any single element is trivial in $\Sthree$. By contrast, the three integers $\{\kone, \ktwo, \kthree\} = \{9, 19, 168\}$ are three distinct values with no non-trivial symmetry under permutation: no non-identity element of $\Sthree$ fixes the ordered triple $(\kone, \ktwo, \kthree)$ as a set (verified in the script, check NEG-b). Hence any bijection $(T_2) \to (T_3)$ must break the $\Sthree$-symmetry of $(T_3)$, and no such bijection is $\Sfour$-equivariant.
\end{proof}

\begin{remark}
This negative result is scientifically informative. It clarifies that $(T_3) = \Vfour\setminus\{e\}$ is the \emph{canonical} carrier of the $\Sthree$-action and the $A_2$ root system, while $(T_2)$ labels the same three objects from the perspective of the leakage dynamics (D51--D52). The canonical bijection between $(T_2)$ and $(T_3)$ is the index-ordering map $\kone \leftrightarrow v_0$, $\ktwo \leftrightarrow v_1$, $\kthree \leftrightarrow v_2$, which is motivated by the quintuplet structure (each index encodes a specific coupling-mode count involving $\nuu$, $\ndd$, and $\Re$) but is not forced by $\Sfour$-equivariance. The relationship between $(T_1)$, $(T_2)$, and $(T_3)$ remains OP-D58-2 in the corpus.
\end{remark}

% ============================================================
\section{Computational verification}
\label{sec:verification}
% ============================================================

The complete derivation of D59 is accompanied by the verification script \texttt{PDL\_D59\_script1.py}~\cite{PDL_D59_script1}, deposited separately on Zenodo. It provides exhaustive computational verification of every lemma using exact integer arithmetic and the Python standard library only (\texttt{fractions}, \texttt{itertools}). No floating-point is introduced.

\begin{table}[H]
\centering
\caption{Items verified by \texttt{PDL\_D59\_script1.py}.}
\label{tab:verification}
\renewcommand{\arraystretch}{1.25}
\small
\begin{tabular}{p{5.0cm} p{6.5cm} l}
\toprule
\textbf{Verification} & \textbf{Method} & \textbf{Status} \\
\midrule
Quintuplet and leakage indices & Direct integer arithmetic & PASSED \\
L1a: $\Sfour$ stabilises $\Vfour\setminus\{e\}$ & Exhaustive conjugation, 24 elements & PASSED \\
L1b: image has 6 elements & Enumeration & PASSED \\
L1c: kernel $= \Vfour$ & Identification, 4 elements & PASSED \\
L2a--e: character decomposition & Fixed-point count, 24 elements & PASSED \\
L3a: $(1,1,1)$ fixed by $\Sthree$ & All 6 actions & PASSED \\
L3b--c: $W$ preserved by $\Sthree$ & Basis vectors, 6 actions & PASSED \\
L3d: 6 distinct $2\times 2$ matrices on $W$ & Explicit computation & PASSED \\
L3e: $|\det| = 1$ for all matrices & Integer determinant & PASSED \\
L4a: norms$^2 = 2$ & Integer dot products & PASSED \\
L4b: $\Sthree$ permutes the 6 roots & Exhaustive & PASSED \\
L4c: transitive orbit & Orbit computation & PASSED \\
L4d: root pairing & Set membership & PASSED \\
L4e--f: Cartan matrix $= A_2$ & Inner products & PASSED \\
L5a--g: orientation from $\Dn > 0$ & Integer inequalities & PASSED \\
NEG-a: $K_i$ distinct & Trivial & PASSED \\
NEG-b: no $\Sthree$ symmetry on $(T_2)$ & Exhaustive, 5 non-trivial perms & PASSED \\
\midrule
\textbf{All verifications} & & \textbf{PASSED} \\
\bottomrule
\end{tabular}
\end{table}

In accordance with the \PDL\ verrouillage protocol, no claim of D59 was incorporated into this document before independent execution of the verification script in Google Colab. The five lemmas and the negative result were verrouillés before LaTeX drafting began.

% ============================================================
\section{Epistemic status}
\label{sec:epistemic}
% ============================================================

\begin{table}[H]
\centering
\caption{Complete epistemic stratification of D59 claims.}
\label{tab:epistemic}
\renewcommand{\arraystretch}{1.3}
\begin{tabular}{p{6.5cm} p{3.0cm} p{4.5cm}}
\toprule
\textbf{Claim} & \textbf{Status} & \textbf{Source} \\
\midrule
$\Dn = 4$ unique (D47) & Theorem (C1--C4) & D47 \\
$\kone = 9$, $\ktwo = 19$, $\kthree = 168$ & Theorem (C1--C4) & D51 \\
$\kone + \ktwo = \ndd$ & Theorem (exact) & D51 \\
$\kthree/(\kone+\ktwo) = \Re = 6$ & Theorem (exact) & D59, L5f \\
$\Sthree$ acts faithfully on $\Vfour\setminus\{e\}$ (L1) & Theorem (C1--C4) & D58 L2, D59 \\
$\chi_\pi = \chi_{\mathbf{1}} + \chi_{\mathbf{2}}$ (L2) & Theorem (group theory) & D59, L2 \\
$W = \{a+b+c=0\}$ is $\Sthree$-invariant (L3) & Theorem (C1--C4) & D59, L3 \\
$A_2$ root system in $W$ (L4) & Theorem (C1--C4) & D59, L4 \\
Cartan matrix $= \bigl[\begin{smallmatrix}2&-1\\-1&2\end{smallmatrix}\bigr]$ & Theorem (classical) & D59, L4 \\
$\Dn > 0$ selects $\mathbf{3}$ over $\bar{\mathbf{3}}$ (L5) & Theorem (C1--C4) & D59, L5 from D47 \\
\textbf{$\SUthreec$ acts in $\mathbf{3}$ on $W$ (main)} & \textbf{Theorem} (classical) & \textbf{D59, Theorem} \\
$(T_2) \not\leftrightarrow (T_3)$ equivariantly & Theorem (verified) & D59, NEG \\
\midrule
$(T_1)$, $(T_2)$, $(T_3)$ canonical relations & Open problem & OP-D58-2 \\
Fermion mass spectrum from C1--C4 & Open problem & OP-D59-1 \\
Three colour generations from C1--C4 & Open problem & OP-OFN-1 \\
\bottomrule
\end{tabular}
\end{table}

The label ``Theorem (classical)'' indicates that the proof additionally invokes the classical Cartan--Killing--Weyl classification and the representation theory of compact Lie groups, taken as established mathematical apparatus.

% ============================================================
\section{Open problems}
\label{sec:openproblems}
% ============================================================

\begin{resolvedproblem}[OP-D58-1 --- Physical identification of $\SUthreec$, resolved by D59]
\label{rp:D58-1}
\textcolor{resolvedgreen}{Resolved by D59.} The fundamental representation $\mathbf{3}$ of $\SUthree$ is constructed on the trace-zero plane $W \subset \mathbb{C}^3$, whose axes are labelled by $\Vfour \setminus \{e\}$. The orientation selecting $\mathbf{3}$ over $\bar{\mathbf{3}}$ is forced by $\Dn = 4 > 0$ (D47).
\end{resolvedproblem}

\begin{resolvedproblem}[OP-D58-3 --- Fundamental representation $\mathbf{3}$ from C1--C4, resolved by D59]
\label{rp:D58-3}
\textcolor{resolvedgreen}{Resolved by D59.} The derivation of $\mathbf{3}$ from C1--C4 is the content of Theorem~\ref{thm:D59_main}.
\end{resolvedproblem}

\begin{openproblem}[OP-D58-2 --- The three triplets of $\Kfour$]
\label{op:D58-2}
Establish formally the relationship between $(T_1)$, $(T_2)$, and $(T_3)$ (Definition~\ref{def:three_triplets}). The negative result of Section~\ref{sec:negative} shows that no $\Sfour$-equivariant bijection $(T_2) \leftrightarrow (T_3)$ exists; the precise nature of their relationship (via the quintuplet) remains to be formalised. \textbf{Priority:} medium.
\end{openproblem}

\begin{openproblem}[OP-D59-1 --- Quark mass spectrum from C1--C4]
\label{op:D59-1}
The present document establishes the carrier space $W$ and the orientation $\mathbf{3}$ for the colour gauge group. It does not provide the quark masses $m_u$, $m_d$, $m_s$, etc. The single irreducible external parameter $\Dm_{\mathrm{iso}} = m_d - m_u$ enters PDL at the QCD interface (D31); the derivation of individual quark masses from C1--C4 is a distinct and harder problem. \textbf{Priority:} medium.
\end{openproblem}

\begin{openproblem}[OP-OFN-1 --- Three leakage cycles and three fermion generations]
\label{op:OFN-1}
The identification of the three leakage cycles with prime exponents $(23, 67, 997)$ with the three generations of fermions, suggested by the common combinatorial origin $\Afour/\Vfour \cong \Zthree$ (D57--D58), remains open. The present document provides additional evidence for this identification (the same triplet $(T_2)$ labels the canonical bijection with the axes of $W$), but a formal derivation requires a mechanism for the mass hierarchy between generations. \textbf{Priority:} high. \textbf{Entry point:} N01, D57, D58, D59.
\end{openproblem}

% ============================================================
\section{Conclusion}
\label{sec:conclusion}
% ============================================================

\subsection{What D59 achieves}

Document D59 completes the programme initiated in D46 and continued in D57--D58 by providing the first explicit construction, from the \PDL\ axioms C1--C4, of a three-dimensional carrier space on which $\SUthree$ acts as the colour gauge group in the fundamental representation $\mathbf{3}$. The main results are:
\begin{enumerate}[leftmargin=2em]
\item The trace-zero plane $W = \{a+b+c=0\} \subset \mathbb{C}^3$ is the canonical carrier of the $\mathbf{3}$ representation, with axes labelled by $\Vfour \setminus \{e\}$ as forced by C1--C4 via $\Kfour$.
\item The $A_2$ root system is realised in $W$ with the Cartan matrix $\bigl[\begin{smallmatrix}2&-1\\-1&2\end{smallmatrix}\bigr]$, confirming at the geometric level what D58 established algebraically.
\item The isospin asymmetry $\Dn = 4 > 0$ (unconditional theorem of C4, D47) selects the matter representation $\mathbf{3}$ over $\bar{\mathbf{3}}$ without additional input.
\item The ratio $\kthree/(\kone+\ktwo) = \Re = 6$ (exact integer) ties the leakage-cycle structure to the relational budget of $\Kfour$, establishing an internal consistency between the cosmological (D51--D52) and gauge (D58--D59) sectors of the corpus.
\end{enumerate}

\subsection{Position within the PDL programme}

The chain D46 + D57 + D58 + D59 now provides:
\begin{itemize}[leftmargin=2em]
\item $\Uone$: U(1) gauge structure and Hopf fibration of $\Kfour$ (D46);
\item $\SUtwo$: SU(2) gauge structure, Weinberg angle $\sin^2\theta_W = 1/4$ at tree level (D57);
\item $\SUthree$: SU(3) gauge structure, algebraic derivation (D58);
\item $\SUthreec$ in $\mathbf{3}$: the colour gauge group with its fundamental representation, matter orientation identified (D59).
\end{itemize}
This constitutes the most complete derivation, from purely combinatorial axioms, of both the gauge group structure and a fundamental gauge representation currently available.

\subsection{What remains open}

D59 does not derive the Standard Model. The phenomenological content --- quark masses, mixing angles, three fermion generations as dynamical replicas --- remains open. The carrier space $W$ is established; the full fermionic content requires additional structure. The programme advances slowly but with certainty.

% ============================================================
\section*{Acknowledgements}

The author thanks Oleg~I.~Evdokimov for the identification of $\Afour/\Vfour \cong \Zthree$ in the joint document N01~\cite{PDL_N01}, which seeded the derivation chain of D57--D59.

% ============================================================
\bibliographystyle{unsrt}
\bibliography{D59_references}

\end{document}
